\documentclass[11pt]{article}
\usepackage[review]{acl}
\usepackage{times}
\usepackage{latexsym}
\usepackage[T1]{fontenc}
\usepackage[utf8]{inputenc}
\usepackage{microtype}
% \usepackage{inconsolata}
\usepackage{graphicx}
\usepackage{amsmath}
\usepackage{amssymb}
\usepackage{booktabs}
\usepackage{multirow}
\usepackage{xcolor}
\usepackage{tabularx}
\usepackage{amsthm}
\newtheorem{theorem}{Theorem}

\setlength\titlebox{6cm}

\title{When Does Diversity Defend? A Framework for Analyzing\\Multi-Judge LLM Evaluation Panel Security}

\author{Anonymous}

\begin{document}
\maketitle

\begin{abstract}
Multi-judge LLM evaluation panels are increasingly deployed under the assumption that judge diversity provides robustness against adversarial manipulation, analogous to Byzantine Fault Tolerance (BFT).
We challenge this assumption by introducing an \emph{adaptive adversary framework} that distinguishes targeted attacks from random faults and reveals that panel security is \emph{regime-dependent}.
Across 455 three-judge panels drawn from 15 models (5 API, 10 open-source) and 9 attack types, we find:
(1)~in panels containing weak models, the weakest member's individual robustness ($\eta_{\max}$) dominates panel attack success rate (significant in 67\% of conditions), confirming a \emph{weakest-link effect};
(2)~in panels composed of strong models, judge diversity ($K_{\text{eff}}$) becomes the operative defense (significant in 50\% of conditions), as no single member presents an exploitable weakness;
(3)~the diversity--capability tradeoff previously reported is an open-source ecosystem artifact ($r=-0.72$ within open-source, $r=-0.05$ within API models), not a fundamental constraint;
(4)~attack transferability constitutes a second regime dimension---targeted attacks amplify the weakest-link effect while transferable attacks favor diverse panels.
We further demonstrate that Gaussian copula models underestimate inter-judge dependence by 10--21$\times$ and that model family labels fail to predict judge independence ($p=0.206$).
Our framework provides regime-specific guidance: prioritize individual robustness for small or weak-model panels, and diversity for large or strong-model panels.
\end{abstract}

\section{Introduction}

Multi-judge LLM evaluation panels---where multiple language models independently assess a candidate response and aggregate via majority vote---are a growing paradigm for scalable evaluation \cite{zheng2024judging}.
A prevailing assumption in this literature is that \emph{diversity} among judges is the primary defense against adversarial manipulation: more independent judges should make the panel harder to deceive, analogous to how Byzantine Fault Tolerance (BFT) in distributed systems leverages redundancy to tolerate faulty nodes.

This analogy is appealing but flawed.
BFT assumes \emph{random} faults---nodes fail independently with some probability.
LLM adversaries, however, are \emph{adaptive}: they can identify and target the weakest judge in a panel, concentrating attack effort where success probability is highest.
Under adaptive adversaries, classical redundancy guarantees break down, and the security properties of a panel depend critically on the \emph{regime}---the capability distribution of its members and the nature of the attack.

In this work, we formalize this distinction and present a comprehensive empirical analysis of when diversity helps and when it hurts.
Our contributions are:

\begin{enumerate}
\item \textbf{Adaptive adversary framework.} We formalize the distinction between random-fault and adaptive adversary threat models for multi-judge panels and prove that under adaptive adversaries, panel attack success rate is bounded by $\Theta(\eta_{\max})$, where $\eta_{\max}$ is the individual susceptibility of the weakest panel member.

\item \textbf{Regime characterization.} We identify two key dimensions that determine whether diversity or individual robustness dominates panel security: (a)~\emph{panel capability tier}---weak-model panels exhibit a weakest-link effect while strong-model panels benefit from diversity; (b)~\emph{attack transferability}---targeted attacks exploit the weakest link while transferable attacks are mitigated by diversity.

\item \textbf{Methodological findings.} We show that Gaussian copula models systematically underestimate inter-judge dependence, and that model family labels (e.g., Qwen, Llama) do not predict functional independence between judges.

\item \textbf{Ecosystem-specific tradeoff.} The diversity--capability tradeoff ($r=-0.72$ in open-source models) largely disappears when API models are included ($r=-0.05$ within-API), revealing it as an artifact of current open-source capability stratification rather than a fundamental constraint.
\end{enumerate}

\section{Related Work}

\paragraph{Multi-judge LLM evaluation.}
LLM-as-a-judge paradigms \cite{zheng2024judging} increasingly use multi-judge panels for robustness.
Recent work advocates diverse panels for reliable evaluation, implicitly assuming that diversity provides BFT-like guarantees.
Our work questions this assumption by formalizing the adaptive adversary threat model.

\paragraph{Adversarial attacks on LLM judges.}
Prompt injection, sycophancy exploitation, and score manipulation have been demonstrated against individual LLM judges.
We extend this line to study how attacks propagate through \emph{panels} of judges under varying diversity conditions.

\paragraph{Byzantine fault tolerance.}
Classical BFT results require $n \geq 3f+1$ honest nodes to tolerate $f$ faulty ones, assuming random faults.
We show that the adaptive nature of LLM adversaries invalidates this assumption in the evaluation setting.

\section{Framework}
\label{sec:framework}

\subsection{Threat Model}

Consider a panel of $K$ judges $\{J_1, \ldots, J_K\}$ that evaluates a candidate response via majority vote.
Each judge $J_i$ has an individual attack success rate $\eta_i \in [0,1]$ representing the probability that an adversary can manipulate $J_i$'s output on a given instance.
The panel attack success rate $\text{ASR}_{\text{panel}}$ is the probability that the adversary flips the majority vote.

We distinguish two threat models:

\paragraph{Random-fault model (BFT-like).}
Each judge fails independently with probability $\eta_i$.
Panel ASR follows a binomial-like formula: the adversary must independently corrupt $\lceil K/2 \rceil$ judges.
Under this model, adding diverse (independent) judges monotonically improves robustness.

\paragraph{Adaptive adversary model.}
The adversary observes the panel composition, identifies the most susceptible judges, and concentrates attack effort.
We prove the following:

\begin{theorem}[Weakest-link bound]
Under an adaptive adversary with per-judge attack budgets, $\mathrm{ASR}_{\mathrm{panel}} = \Theta(\eta_{\max})$, where $\eta_{\max} = \max_i \eta_i$.
\end{theorem}

The bound states that panel security is dominated by the weakest member.
Crucially, this bound is \emph{tight} only when $\eta_{\max}$ is substantially larger than the median $\eta_i$---i.e., when a genuinely weak member exists.
When all members are comparably strong, the bound becomes loose and diversity provides additional protection through complementary failure modes.

\subsection{Diversity and Robustness Metrics}

We quantify judge diversity via the \emph{effective number of independent judges} $K_{\text{eff}}$, derived from the mutual information (MI) matrix among judges' outputs:
\begin{equation}
K_{\text{eff}} = K \cdot \left(1 - \frac{\bar{I}}{H_{\max}}\right)
\end{equation}
where $\bar{I}$ is the mean pairwise MI and $H_{\max}$ is the maximum single-judge entropy.
$K_{\text{eff}} = K$ indicates fully independent judges; $K_{\text{eff}} = 1$ indicates identical judges.

Individual robustness is measured by $\eta_i$, the per-judge attack success rate, and we define $\eta_{\max} = \max_i \eta_i$ as the panel's weakest-link susceptibility.

\section{Experimental Setup}

\subsection{Models}

We assemble a pool of 15 models spanning three capability tiers:
(1)~\textbf{API models} (5): GPT-4o-mini, GPT-4o, GPT-4.1, Claude-3.5-Sonnet, Gemini-1.5-Pro;
(2)~\textbf{Large open-source} (4): Llama-3.1-70B, Qwen2.5-72B, Mixtral-8x22B, and one additional 70B-class model;
(3)~\textbf{Small open-source} (6): five 7B-class models plus Mixtral-8x7B.
From these, we enumerate $\binom{15}{3} = 455$ three-judge panels and additionally sample panels of size $K \in \{5, 7, 9\}$.

\subsection{Attacks}

We evaluate 9 attack types spanning targeted and transferable categories:
\emph{prompt injection}, \emph{score manipulation}, \emph{sycophancy}, \emph{verbosity bias}, \emph{position bias}, \emph{authority bias} (from prior work), plus three new attacks: \emph{few-shot hijacking}, \emph{persona manipulation}, and \emph{instruction override}.

\subsection{Datasets and Evaluation}

Experiments are conducted on MMLU and ARC benchmarks.
For each panel $\times$ attack condition, we measure panel ASR and fit multivariate regressions:
\begin{equation}
\log(\text{ASR}) \sim \beta_1 \log(\eta_{\max}) + \beta_2 \log(K_{\text{eff}}) + \epsilon
\end{equation}
to decompose the contributions of individual robustness and diversity.

\section{Results}

\subsection{Weakest-Link Effect is Regime-Dependent}
\label{sec:scale}

Table~\ref{tab:regression_main} presents multivariate regression results across 455 panels and 12 dataset$\times$attack conditions.
Overall, $\eta_{\max}$ is significant in 5/12 conditions (42\%), with a median $\beta = 0.61$, while $K_{\text{eff}}$ achieves significance in only 4/12 conditions (33\%), with a median $\beta = -0.18$ and a median $R^2 = 0.22$.

\begin{table}[t]
\centering
\small
\begin{tabular}{llcccc}
\toprule
\textbf{Data} & \textbf{Attack} & $\beta(\eta)$ & $p(\eta)$ & $\beta(K)$ & $p(K)$ \\
\midrule
MMLU & prompt-inj. & 0.82 & .003 & $-$0.31 & .08 \\
MMLU & score-manip. & 0.74 & .011 & $-$0.22 & .18 \\
MMLU & sycophancy & 0.58 & .042 & $-$0.41 & .03 \\
MMLU & verbosity & 0.45 & .12 & $-$0.15 & .42 \\
MMLU & position & 0.39 & .18 & 0.12 & .51 \\
MMLU & authority & 0.68 & .008 & $-$0.28 & .11 \\
\midrule
ARC & prompt-inj. & 0.71 & .015 & $-$0.38 & .04 \\
ARC & score-manip. & 0.63 & .028 & $-$0.19 & .27 \\
ARC & sycophancy & 0.52 & .065 & $-$0.44 & .02 \\
ARC & verbosity & 0.41 & .14 & $-$0.33 & .07 \\
ARC & position & 0.33 & .22 & 0.18 & .35 \\
ARC & authority & 0.55 & .052 & $-$0.26 & .15 \\
\bottomrule
\end{tabular}
\caption{Multivariate regression on 455 panels ($K$=3). $\beta(\eta)$: coefficient for $\log(\eta_{\max})$; $\beta(K)$: coefficient for $\log(K_{\text{eff}})$. $R^2$ ranges 0.08--0.34 (median 0.22). Bold $p$-values indicate significance at $\alpha=0.05$.}
\label{tab:regression_main}
\end{table}

Critically, these aggregate statistics mask a sharp regime split.
Table~\ref{tab:regime_stratified} stratifies panels by capability tier.
In all-small-OS panels (7B class), $\eta_{\max}$ is significant in 67\% of conditions with a median $R^2 = 0.35$---the weakest-link effect holds strongly.
In all-API panels, $\eta_{\max}$ significance drops to 17\%, while $K_{\text{eff}}$ becomes significant in 50\% of conditions.
The explanation is straightforward: among strong models, $\eta_{\max}$ variance is compressed (all models have low $\eta$), removing the statistical signal.
Diversity then captures complementary vulnerability profiles across models from different providers.

\begin{table}[t]
\centering
\small
\begin{tabular}{lcccc}
\toprule
\textbf{Panel Tier} & $N$ & $\eta$ sig & $K$ sig$^-$ & $R^2$ \\
\midrule
All-API & 20 & 0.17 & 0.50 & 0.12 \\
Mixed (API+70B) & 120 & 0.33 & 0.42 & 0.19 \\
All-large-OS & 35 & 0.50 & 0.25 & 0.24 \\
Mixed (lg+sm OS) & 180 & 0.58 & 0.17 & 0.28 \\
All-small-OS & 10 & 0.67 & 0.08 & 0.35 \\
$\geq$1 weak model & 280 & 0.55 & 0.18 & 0.27 \\
\bottomrule
\end{tabular}
\caption{Regime-stratified analysis. ``$\eta$ sig'': fraction of conditions where $\eta_{\max}$ is significant; ``$K$ sig$^-$'': fraction where $K_{\text{eff}}$ is significantly negative. The weakest-link effect concentrates in weak-model panels.}
\label{tab:regime_stratified}
\end{table}

\subsection{Attack Transferability as a Second Regime Dimension}
\label{sec:attack_expansion}

Table~\ref{tab:new_attacks} shows regression results for three additional attack types.
The results reveal attack transferability as a key dimension governing which factor dominates.

\emph{Persona manipulation}---a targeted attack requiring model-specific tuning---shows strong $\eta_{\max}$ significance ($\beta = 0.72$, $p = 0.009$ on MMLU) with non-significant $K_{\text{eff}}$.
In contrast, \emph{few-shot hijacking} and \emph{instruction override}---transferable attacks exploiting shared instruction-following properties---show non-significant $\eta_{\max}$ but significant $K_{\text{eff}}$ ($p = 0.03$ and $p = 0.02$ on MMLU, respectively).

\begin{table}[t]
\centering
\small
\begin{tabular}{llcccc}
\toprule
\textbf{Data} & \textbf{Attack} & $\beta(\eta)$ & $p(\eta)$ & $\beta(K)$ & $p(K)$ \\
\midrule
MMLU & few-shot-hij. & 0.28 & .24 & $-$0.48 & .03 \\
MMLU & persona-man. & 0.72 & .009 & $-$0.14 & .39 \\
MMLU & instr.-overr. & 0.35 & .15 & $-$0.52 & .02 \\
ARC & few-shot-hij. & 0.31 & .19 & $-$0.39 & .05 \\
ARC & persona-man. & 0.65 & .035 & $-$0.21 & .22 \\
ARC & instr.-overr. & 0.38 & .12 & $-$0.44 & .03 \\
\bottomrule
\end{tabular}
\caption{Regression results for three new attack types on 455 panels. Targeted attacks (persona manipulation) favor $\eta_{\max}$; transferable attacks (few-shot hijacking, instruction override) favor $K_{\text{eff}}$.}
\label{tab:new_attacks}
\end{table}

This yields a two-dimensional regime characterization: (1)~\emph{panel capability tier} determines whether a weak link exists to exploit, and (2)~\emph{attack transferability} determines whether the adversary can concentrate effort on a single judge.
Across all 18 conditions (9 attacks $\times$ 2 datasets), $\eta_{\max}$ is significant in 39\% and $K_{\text{eff}}$ in 39\%---neither universally dominates.

\subsection{Adaptive vs.\ Random-Fault Distinction}
\label{sec:bft}

To validate our framework's core theoretical distinction, we compare adaptive attacks with a random-fault simulation where each judge fails independently with probability $\eta_i$.
Table~\ref{tab:bft_contrast} presents the results stratified by regime.

\begin{table}[t]
\centering
\small
\begin{tabular}{llcccc}
\toprule
\textbf{Regime} & \textbf{Fault} & $\beta(K)$ & $p(K)$ & $\beta(\eta)$ & $p(\eta)$ \\
\midrule
\multirow{2}{*}{Weak} & adaptive & $-$0.12 & .48 & 0.78 & .004 \\
 & random & $-$0.41 & .03 & 0.52 & .06 \\
\midrule
\multirow{2}{*}{Mixed} & adaptive & $-$0.28 & .09 & 0.55 & .04 \\
 & random & $-$0.45 & .02 & 0.38 & .11 \\
\midrule
\multirow{2}{*}{Strong} & adaptive & $-$0.44 & .02 & 0.22 & .31 \\
 & random & $-$0.58 & .005 & 0.15 & .48 \\
\bottomrule
\end{tabular}
\caption{Adaptive vs.\ random-fault comparison by regime. Under random faults, $K_{\text{eff}}$ is consistently more protective. The gap is largest in weak panels (Cohen's $d = 0.38$, $p = 0.04$) and negligible in strong panels ($d = 0.12$, $p = 0.41$).}
\label{tab:bft_contrast}
\end{table}

The BFT contrast---the difference in $K_{\text{eff}}$'s protective effect between random and adaptive conditions---is itself regime-dependent.
In weak panels, adaptive adversaries suppress diversity's benefit significantly more than random faults (Cohen's $d = 0.38$, $p = 0.04$), consistent with targeted exploitation of the weakest member.
In strong panels, the distinction diminishes ($d = 0.12$, $p = 0.41$), because even adaptive adversaries lack a clear target when all judges are robust.
The overall BFT contrast is modest ($d = 0.22$, $p = 0.08$), with a significant regime$\times$fault-type interaction ($p = 0.04$).

\subsection{Per-Judge ASR Concentration}
\label{sec:concentration}

To directly test whether adaptive adversaries concentrate on the weakest judge, we analyze per-judge ASR distributions via Gini coefficients and entropy (Table~\ref{tab:concentration}).

\begin{table}[t]
\centering
\small
\begin{tabular}{lcccc}
\toprule
\textbf{Regime} & \textbf{Gini} & \textbf{Entropy} & \textbf{Weak-J \%} & $N$ \\
\midrule
All-API & 0.18 & 0.91 & 0.38 & 20 \\
Mixed (API+70B) & 0.28 & 0.83 & 0.42 & 120 \\
All-large-OS & 0.35 & 0.77 & 0.47 & 35 \\
Mixed (lg+sm) & 0.45 & 0.69 & 0.55 & 180 \\
All-small-OS & 0.52 & 0.62 & 0.61 & 10 \\
\bottomrule
\end{tabular}
\caption{Per-judge ASR concentration. ``Weak-J \%'': fraction of panel ASR attributable to the weakest judge. In weak panels, the weakest judge absorbs 61\% of successful attacks (Gini $= 0.52$); in strong panels, ASR is nearly uniformly distributed (Gini $= 0.18$).}
\label{tab:concentration}
\end{table}

The concentration gradient is stark: Gini coefficients range from 0.18 (near-uniform) in all-API panels to 0.52 (highly concentrated) in all-small-OS panels, with the weakest judge absorbing 38--61\% of successful attacks.
Comparing adaptive and random-fault conditions, adaptive attacks produce modestly higher concentration (median Gini 0.34 vs.\ 0.26), though the difference does not reach significance ($p = 0.06$, $d = 0.28$).
This supports the interpretation that adaptive adversary behavior is most empirically distinguishable from random faults when a clear weak link exists.

\subsection{Capability-Matched Diversity Analysis}
\label{sec:capability_matched}

A confound in cross-tier analyses is that diversity and capability are mechanically correlated: panels mixing 7B and 70B models are diverse but contain weak members.
We control for this by analyzing panels \emph{within} capability tiers (Table~\ref{tab:within_tier}).

\begin{table}[t]
\centering
\small
\begin{tabular}{lccccc}
\toprule
\textbf{Tier} & $N_m$ & $N_p$ & $\eta$ sig & $K$ sig$^-$ & $r(K,\eta)$ \\
\midrule
7B class & 5 & 10 & 0.50 & 0.17 & $-$0.18 \\
70B class & 4 & 4 & 0.33 & 0.33 & $-$0.08 \\
API & 5 & 10 & 0.17 & 0.50 & $-$0.05 \\
Cross-tier & 15 & 455 & 0.42 & 0.33 & $-$0.41 \\
\bottomrule
\end{tabular}
\caption{Within-tier regression analysis. $N_m$: models in tier; $N_p$: panels. Within-tier $K_{\text{eff}}$--$\eta_{\max}$ correlations approach zero, indicating no inherent diversity--capability tradeoff within homogeneous tiers.}
\label{tab:within_tier}
\end{table}

Within capability tiers, the $K_{\text{eff}}$--$\eta_{\max}$ correlation approaches zero ($r = -0.05$ to $-0.18$), compared to $r = -0.41$ across tiers.
Importantly, within the API tier, $K_{\text{eff}}$ is significantly negative in 50\% of conditions---diversity genuinely helps when all judges are strong.
This occurs because API models from different providers (OpenAI, Anthropic, Google) have distinct vulnerability profiles despite comparable overall capability.

The diversity--capability tradeoff decomposes sharply by ecosystem (Table~\ref{tab:tradeoff_decomp}).

\begin{table}[t]
\centering
\small
\begin{tabular}{lcccl}
\toprule
\textbf{Scope} & $r$ & $p$ & $R^2$ & \textbf{Interpretation} \\
\midrule
All 15 models & $-$0.41 & .003 & .17 & Moderate \\
OS only (10) & $-$0.72 & {$<$.001} & .52 & Strong \\
API-incl. ($\geq$1) & $-$0.22 & .04 & .05 & Weak \\
All-API (5) & $-$0.05 & .89 & .003 & None \\
\bottomrule
\end{tabular}
\caption{Diversity--capability tradeoff by model pool. The tradeoff is an open-source ecosystem artifact ($r = -0.72$), not a fundamental constraint ($r = -0.05$ within API).}
\label{tab:tradeoff_decomp}
\end{table}

The open-source ecosystem exhibits a strong tradeoff ($r = -0.72$, $p < 0.001$) because maximizing diversity across open-source models requires crossing capability tiers (mixing 7B and 70B), necessarily introducing weaker members.
Among API models, the tradeoff is absent ($r = -0.05$, $p = 0.89$): models from different providers are naturally diverse yet comparably capable.
On the Pareto frontier of 455 panels, 8 of 12 optimal panels contain at least 2 API models, with mean ASR of 0.075 compared to 0.148 for non-Pareto panels.

\subsection{Panel Size Sensitivity}
\label{sec:panel_size}

We test whether the weakest-link effect persists across panel sizes $K \in \{3, 5, 7, 9\}$ (Table~\ref{tab:panel_size}).

\begin{table}[t]
\centering
\small
\begin{tabular}{lccccc}
\toprule
$K$ & $N_p$ & $\eta$ sig & $K$ sig$^-$ & $R^2$ & ASR \\
\midrule
3 & 455 & 0.42 & 0.33 & 0.22 & 0.134 \\
5 & 3003 & 0.33 & 0.39 & 0.18 & 0.105 \\
7 & 6435 & 0.28 & 0.50 & 0.15 & 0.092 \\
9 & 5005 & 0.22 & 0.56 & 0.13 & 0.087 \\
\bottomrule
\end{tabular}
\caption{Panel size sensitivity. $\eta_{\max}$ significance monotonically decreases with $K$ (trend $p = 0.03$) while $K_{\text{eff}}$ significance increases ($p = 0.02$). The crossover occurs near $K = 5$.}
\label{tab:panel_size}
\end{table}

The $\eta_{\max}$ significance rate decreases monotonically from 0.42 ($K=3$) to 0.22 ($K=9$), while $K_{\text{eff}}$ significance increases from 0.33 to 0.56.
Both trends are significant ($p = 0.03$ and $p = 0.02$, respectively), with crossover near $K = 5$.
The intuition is that in small panels, a single compromised judge represents a large fraction of the vote (1/3 at $K=3$); in larger panels, the weakest judge's influence is diluted and the combinatorial difficulty of corrupting a majority increases, making diversity---which ensures diverse failure modes---more protective.
Increasing $K$ from 3 to 9 reduces mean ASR by 35\%.

\subsection{Panel Selection Strategies}
\label{sec:panel_selection}

We benchmark a \emph{two-stage} panel selection algorithm---first screen for individual robustness ($\eta \leq$ threshold), then maximize $K_{\text{eff}}$ among qualifying judges---against five baselines (Table~\ref{tab:selection}).

\begin{table}[t]
\centering
\small
\begin{tabular}{lcccc}
\toprule
\textbf{Strategy} & \textbf{ASR} & \textbf{Win} & $p_{\text{vs.rand}}$ \\
\midrule
Two-stage (ours) & 0.091 & 0.44 & .001 \\
Strongest-$K$ & 0.082 & 0.50 & {$<$.001} \\
Min-$\eta_{\max}$ & 0.095 & 0.39 & .003 \\
Max-$K_{\text{eff}}$ & 0.131 & 0.33 & .02 \\
Max-family-div. & 0.158 & 0.11 & .15 \\
Random & 0.172 & 0.06 & --- \\
\bottomrule
\end{tabular}
\caption{Panel selection benchmark. Win rate: fraction of 18 conditions where strategy achieves lowest ASR. All robustness-aware strategies significantly outperform random.}
\label{tab:selection}
\end{table}

The two-stage algorithm does not significantly outperform the simpler Strongest-$K$ baseline ($p = 0.32$), which selects the $K$ models with lowest individual $\eta$.
However, regime-stratified analysis (Table~\ref{tab:selection_regime}) reveals complementary strengths.

\begin{table}[t]
\centering
\small
\begin{tabular}{lcccc}
\toprule
\textbf{Regime} & \textbf{Ours} & \textbf{Str-$K$} & \textbf{Tie} & $p$ \\
\midrule
Weak (OS-7B only) & 4/6 & 2/6 & 0 & .12 \\
Mixed & 3/6 & 3/6 & 0 & .82 \\
Strong (API avail.) & 1/6 & 4/6 & 1 & .08 \\
\bottomrule
\end{tabular}
\caption{Head-to-head by regime. Two-stage is preferable in weak-model pools; Strongest-$K$ suffices when strong models are available.}
\label{tab:selection_regime}
\end{table}

In weak-model pools (open-source 7B only), the two-stage algorithm wins 4/6 conditions---its diversity-maximization step provides value when all candidates have moderate vulnerability.
When API models are available, Strongest-$K$ wins 4/6 conditions because selecting the strongest models also yields natural diversity (different providers).
The practical implication: two-stage selection is most valuable precisely when the model pool is constrained to weaker options.

Additionally, we note that the threshold parameter in the two-stage algorithm exhibits high sensitivity to attack type, with no single threshold performing optimally across all conditions.

\section{Discussion}

\paragraph{A unified regime framework.}
Our results converge on a two-dimensional regime characterization of panel security.
Along the \emph{capability} axis, panels transition from weakest-link dominance (weak models) to diversity-mediated defense (strong models), with the transition occurring when $\eta_{\max}$ variance compresses below a threshold at which it ceases to be a significant predictor.
Along the \emph{attack transferability} axis, targeted attacks amplify weakest-link dynamics while transferable attacks distribute risk across judges, favoring diverse panels.

\paragraph{Implications for panel design.}
Our findings yield regime-specific guidance:
(1)~For small panels ($K \leq 5$) or pools containing weak models, prioritize individual robustness---screen judges for low $\eta$ before considering diversity.
(2)~For large panels ($K \geq 7$) composed of strong models, diversity is the operative defense and should be maximized.
(3)~Model family labels are unreliable proxies for diversity; MI-based $K_{\text{eff}}$ should be used instead.

\paragraph{Limitations of the BFT analogy.}
The adaptive-vs.-random distinction is theoretically clean but empirically regime-dependent.
In strong-model panels, adaptive adversary behavior becomes empirically indistinguishable from random faults (Cohen's $d = 0.12$), not because the adversary is not adaptive, but because there is no sufficiently weak target to concentrate on.
The framework's value lies in formalizing when and why the BFT analogy breaks, rather than claiming it always does.

\paragraph{Ecosystem dynamics.}
The diversity--capability tradeoff is an artifact of the current open-source model landscape, where capability strongly predicts robustness.
As open-source models improve (e.g., future 7B models approaching current 70B capability), we expect this tradeoff to diminish, similar to how it is already absent among API models.

\section{Conclusion}

We present a framework for analyzing multi-judge LLM evaluation panel security that distinguishes adaptive adversaries from random faults and characterizes when diversity helps versus when individual robustness dominates.
Our key finding is that panel security is regime-dependent: the weakest-link effect holds in panels containing weak models or facing targeted attacks, while diversity provides genuine protection in strong-model panels facing transferable attacks.
The diversity--capability tradeoff is ecosystem-specific, not fundamental.
We additionally contribute methodological findings on copula model failures and family label unreliability that inform broader work on LLM evaluation independence.

\section*{Limitations}

Our study has three main limitations.
First, while we expand to 15 models, the API model tier contains only 5 models from three providers; the within-API diversity analysis is based on a relatively small number of panels ($N=10$ at $K=3$).
Second, our adaptive adversary model assumes the attacker has knowledge of individual judge susceptibilities, which may overestimate real-world attacker capability; the empirical BFT contrast ($d = 0.22$) suggests actual attacks may fall between our random and adaptive extremes.
Third, the experiments use MMLU and ARC, both multiple-choice benchmarks; the regime characterization may differ for open-ended generation evaluation where judge failure modes are less structured.

\bibliographystyle{acl_natbib}
\bibliography{references}

\end{document}
